\documentclass[11pt,letterpaper]{article}
\usepackage[margin=0.72in]{geometry}
\usepackage[T1]{fontenc}
\usepackage[utf8]{inputenc}
\usepackage{lmodern}
\usepackage{microtype}
\usepackage{xcolor}
\usepackage{enumitem}
\usepackage{tabularx}
\usepackage{hyperref}

\definecolor{accent}{HTML}{00419B}
\definecolor{muted}{HTML}{555555}

\hypersetup{
  colorlinks=true,
  urlcolor=accent,
  linkcolor=accent
}

\setlength{\parindent}{0pt}
\setlength{\parskip}{0.45em}
\setlist[itemize]{leftmargin=1.25em, topsep=0.2em, itemsep=0.15em}

\newcommand{\MauriName}{Maurizio Ungaro}
\newcommand{\MauriTitle}{Staff Scientist, Jefferson Lab}
\newcommand{\MauriEmail}{\href{mailto:ungaro@jlab.org}{ungaro@jlab.org}}
\newcommand{\MauriHomepage}{\href{https://maureeungaro.github.io/home/}{maureeungaro.github.io/home}}
\newcommand{\MauriGitHub}{\href{https://github.com/maureeungaro}{github.com/maureeungaro}}
\newcommand{\MauriScholar}{\href{https://scholar.google.com/citations?user=zkWYILYAAAAJ\&hl=en}{Google Scholar}}
\newcommand{\MauriInspire}{\href{https://inspirehep.net/authors/1322331}{INSPIRE}}
\newcommand{\MauriOrcid}{\href{https://orcid.org/0000-0001-6982-3310}{ORCID 0000-0001-6982-3310}}
\newcommand{\MauriLinkedIn}{\href{https://www.linkedin.com/in/maurizio-ungaro-37992062/}{LinkedIn}}
\newcommand{\MauriJLab}{\href{https://www.jlab.org}{Jefferson Lab}}
\newcommand{\MauriHallB}{\href{https://www.jlab.org/physics/hall-b}{Jefferson Lab Hall-B}}
\newcommand{\MauriRPI}{\href{https://www.rpi.edu}{Rensselaer Polytechnic Institute}}
\newcommand{\MauriUniGe}{\href{https://unige.it/}{Universita degli Studi di Genova}}
\newcommand{\MauriGEMC}{\href{https://gemc.github.io/home/}{GEMC}}
\newcommand{\MauriGeantFourJLab}{\href{https://jeffersonlab.github.io/g4home/}{Geant4 at JLab}}
\newcommand{\MauriClasTwelve}{\href{https://github.com/gemc/clas12Tags}{CLAS12 Simulations}}
\newcommand{\MauriOSG}{\href{https://maureeungaro.github.io/home/osg/osg}{CLAS12 on OSG}}
\newcommand{\MauriLTCC}{\href{https://maureeungaro.github.io/home/ltcc/ltcc}{Low Threshold Cherenkov Counter}}
\newcommand{\MauriSPringEight}{\href{https://www.spring8.or.jp/en/}{SPring-8}}
\newcommand{\MauriHPS}{\href{https://www.jlab.org/physics/hall-b/other}{Heavy Photon Search}}

\newcommand{\DocHeader}[1]{
  {\LARGE\bfseries \MauriName}\hfill {\color{muted}#1}\\
  Staff Scientist, \MauriJLab{}\\
  \MauriEmail{} \quad
  \MauriHomepage{} \quad
  \MauriGitHub{}\\
  \MauriScholar{} \quad
  \MauriInspire{} \quad
  \MauriOrcid{}\\
  \vspace{0.4em}
  \hrule
  \vspace{0.8em}
}

\newcommand{\RoleEntry}[4]{
  \textbf{#1}\hfill {\color{muted}#2}\\
  #3\\
  #4
}

\newcommand{\ProjectEntry}[3]{
  \textbf{#1}: #2\\
  {\small\href{#3}{#3}}
}


\begin{document}
\DocHeader{Application Letter Templates}

\section*{Purpose}
This document preserves reusable application-letter material from older
academic applications. It is intentionally not installed as a public homepage
asset. Use it as source material when drafting targeted applications.

\section*{Research-Focused Opening}
I am a nuclear physicist and scientific software developer at \MauriJLab{} whose
work connects nucleon-structure research, detector systems, and production-scale
simulation infrastructure. I develop and support Geant4/\MauriGEMC{} workflows,
maintain CLAS12 production tools, contribute to detector operation and
calibration, and write documentation that helps users apply complex simulation
tools productively.

\section*{Teaching-Focused Opening}
I am interested in positions where teaching, mentoring, and experimental
research strengthen one another. My teaching experience includes undergraduate
physics laboratories at Christopher Newport University, teaching assistantships
at \MauriRPI{}, and extensive mentoring of students working on data analysis,
detector calibration, simulation, and CLAS/CLAS12 physics at Jefferson Lab.

\section*{Research Program Paragraph}
My research focuses on the internal structure and dynamics of the nucleon,
including meson electro-production, nucleon resonances, and the transition
between hadronic and partonic degrees of freedom. This work connects naturally
to Jefferson Lab programs in exclusive reactions, generalized parton
distributions, quark propagation in nuclei, and dark-sector searches through
\MauriHPS{}.

\section*{Teaching Philosophy Paragraph}
My teaching approach emphasizes active problem solving and physical intuition.
Students learn most effectively when they connect equations to measurements:
what can be observed, how it can be measured, how uncertainties enter, and how
data and simulation test an idea. I use modern computing, data acquisition,
visualization, and simulation tools to make laboratory work concrete and to help
students build confidence as independent problem solvers.

\section*{Closing Paragraph}
I would welcome the opportunity to discuss how my experience in experimental
nuclear physics, scientific software, detector simulation, teaching, and
mentoring can support your program. My background is strongest where research
quality, practical computing, and student development all matter.

\end{document}
